\section{Dependency audit for the monopole coefficient}
\label{supp:hypothesis-audit}
The five hypotheses in the main text describe dependencies of the
coefficient enclosure. The scalar summation alone does not prove its
input lemmas. This register distinguishes inherited computer-assisted
results, analytic implications, conventions and the additional
regularity assumption on a hypothetical nonlinear family. The
accompanying archive is distributed with release \texttt{qball-bic-v0.51}
(see the data availability statement of the main paper).

\subsection{Status of the five dependencies}
\begin{enumerate}
\item\textbf{Common exact background and mode: inherited existence result.}
The defining C2 certificate, \texttt{ZERT-T2-A-L36.json} in
\texttt{archive/proofs/C1-C2/}, encloses one common exact tuple
$z_*=(a_*,c_*,\rho_*,\omega_*)$. The analytic implication is the
linear theorem together with the origin, tail and inclusion lemmas in
Section~\ref{supp:proof-lemmas}. All following enclosures refer to this
same tuple. Neither a distinct fit nor every point in the parameter box
is assumed to be an exact root. The code and arithmetic trust assumption
is inherited; uniqueness of the tuple is unnecessary.

\item\textbf{Regular/Jost identification: inherited Wronski argument.}
At this tuple, the two exact matching Wronskians vanish, so the
normalized decaying Jost vector belongs to the two-dimensional regular
origin span. Section~\ref{supp:matching-positive} supplies this analytic
step. The archived matching-component bounds enclose its coefficients
in that span. Their residual intervals containing zero would not, by
themselves, establish a common solution. The exterior evaluation binds
the C2 box and verifies the rational majorants
\begin{equation}
 |f_*(r)|\leq\frac{97}{8}\frac{e^{-0.495r}}{r},\qquad
 |u(r)|+|v(r)|\leq1.097e^{-0.287r},\qquad r\geq36.
 \label{supp:exact-tails}
\end{equation}
These are the fundamental mode's reduced components and Jost
normalization, not unit-$L^2$ data.

\item\textbf{Source and normalization: definition plus checked binding.}
The exact source is the full quadratic product in the main text.
The fixed nominal unit-$L^2$ mode is divided by
$s_0=2558188928468177/2251799813685248$, and quadratic quantities by
$s_0^2$. This exact rational conversion defines the comparison data;
it does not set the exact mode's $L^2$ norm equal to $s_0$.
The source and response exports use the same nominal functions and
carrier conjugation. Exact-to-nominal errors remain in the budget.
This item is a convention and input-identity check, rather than a
new existence assumption. The archive's normalization and error-identity
lemmas specify it independently of the stored scalar sum.

\item\textbf{Exact outgoing problem: analytic implication with checked inputs.}
The reduced second-harmonic operator has a real symmetric locally
bounded potential and both $q_\pm$ strictly between $5/2$ and $23/5$.
These inequalities are checked over the same exact C2 box. The exterior
evaluation gives, at $R=64$,
\begin{equation}
 \int_{64}^{\infty}\|V_*\|_2\,\dd r<10^{-28},\qquad
 \int_{64}^{\infty}\|F_*\|_2\,\dd r<10^{-32}.
 \label{supp:exterior-masses}
\end{equation}
In particular the contraction condition
$\int_R^\infty\|V_*\|_2/q_{\min}<1/2$ is verified at the exact
frequencies. The construction below includes the exterior source and
proves the unique regular outgoing problem and its adjoint extraction.
Thus this item is not an untested numerical Robin-condition assumption.
It still depends on the inherited exact-mode and tail lemmas and the
correctness of their validated evaluation.

\item\textbf{Quantitative bounds and complete budget: checked inputs and sum.}
The complete-potential enclosure gives
$\int_0^\infty\|V_*\|_2\,\dd r<25/4$. The analytic argument below
therefore supplies $\|Z\|_2<26/5$, $\|Z'\|_2<14$ for the exact
adjoint. The archived component budget bounds source and operator
errors, the actual reconstructed response residual, origin terms,
matching jumps, exterior corrections and coefficient extraction.
Its classification is a satisfied conditional coefficient criterion;
there is no unfilled numerical summand in this $L=0$ budget.
Accepting the historical input chain and its directed-arithmetic
execution remains a trust assumption, rather than a new independent
verification performed for this manuscript revision.
\end{enumerate}

All five items have identified analytical or computational support.
The adjective \emph{conditional} records this inherited proof chain and
its software assumptions; it does not designate an omitted term of the
reported scalar budget. The separate assumption that a hypothetical
nonlinear family is twice continuously differentiable in a fixed
weighted space restricts the family excluded by the continuation
argument. The circular-polarization case and a quantitative quadrupole
nonzero certificate remain outside this monopole conclusion.

\subsection{Exterior reduction and unique outgoing response}
Let $-U''-Q^2U+VU=F$, with $Q=\operatorname{diag}(q_+,q_-)$ and
$q_j\geq q_{\min}>0$. In the computational convention both outgoing
carriers are $e^{iq_jr}$. Define
\begin{equation}
 (Kg)(r)=\int_r^\infty Q^{-1}\sin(Q(s-r))g(s)\,\dd s,\quad
 E(r)=e^{iQr},\quad H=E+KVH,\quad p=KVp-KF.
\end{equation}
Set $A_R=\int_R^\infty\|V\|_2$, $B_R=\int_R^\infty\|F\|_2$,
$\eta=A_R/q_{\min}<1/2$, $\delta=\eta/(1-\eta)$ and
$\gamma_R=1/(1-\delta)$. Neumann inversion on bounded functions gives
\begin{equation}
 \|H\|\leq\frac1{1-\eta},\quad \|H-E\|\leq\delta,\quad
 \|H'-E'\|\leq\frac{A_R}{1-\eta},\quad
 \|p\|\leq\frac{B_R}{q_{\min}(1-\eta)},\quad
 \|p'\|\leq\frac{B_R}{1-\eta}.
\end{equation}
The corrections and their first derivatives vanish at infinity.
Since $E(R)$ is unitary and $\delta<1$, $H(R)$ is invertible with
inverse norm at most $\gamma_R$. Thus $U=HC+p$ has the exact maps
\begin{equation}
 \begin{aligned}
 T=H(R)^{-1},\quad c=-Tp(R),\quad D=H'(R)T,\quad
 &b=p'(R)-Dp(R),\\
 C=TU(R)+c,\qquad &U'(R)=DU(R)+b.
 \end{aligned}
\end{equation}
The source contributes both $b$ and $c$. Dropping it would change the
problem even if the homogeneous maps were exact.
For $d_R=[A_R/(1-\eta)+q_{\max}\delta]\gamma_R$, the bounds are
\begin{equation}
 \|T-E(R)^{-1}\|\leq\delta\gamma_R,\quad
 \|D-iQ\|\leq d_R,\quad
 \|c\|\leq\gamma_R\|p\|,\quad
 \|b\|\leq\|p'\|+(q_{\max}+d_R)\|p\|.
\end{equation}
They follow by inserting $H-E$ into the inverse identity and into
$H'-iQH$; all matrix products keep their displayed order.

For real symmetric $V$, the conserved matrix current is
$(H^\dagger H'-H'^\dagger H)/(2i)=Q$. Consequently
\begin{equation}
 \frac{D-D^\dagger}{2i}=H(R)^{-\dagger}QH(R)^{-1}>0.
\end{equation}
A homogeneous regular solution $U(0)=0$ satisfying $U'(R)=DU(R)$
has real interior Green energy. Its boundary imaginary part is positive
unless $U(R)=0$. Hence both Cauchy data vanish and $U=0$.
The finite regular-to-boundary matrix is therefore invertible, which
proves existence and uniqueness for the inhomogeneous boundary data.
The adjoint problem uses $D^\dagger$ and has negative definite boundary
imaginary part, so the same argument proves its uniqueness.
This reasoning applies to the two open second-harmonic channels.
The fundamental embedded mode has a closed channel and does not satisfy
this positivity premise.

\subsection{Global adjoint weights and the error identity}
Let $A(r)=\int_r^\infty\|V(s)\|_2\,\dd s$ and
$h(r)=e^{A(r)/q_{\min}}$. Ordered Volterra iteration on the whole
half-line gives outgoing and incoming matrices $H_\pm$ with
\begin{equation}
 \|H_\pm(r)\|_2\leq h(r),\qquad
 \|H_\pm'(r)\|_2\leq q_{\max}+q_{\min}(h(r)-1).
\end{equation}
No smallness condition on $A(0)$ is required: the $n$th integral has
bound $(A(r)/q_{\min})^n/n!$. The conserved sesquilinear Wronskians are
$W(H_\pm,H_\pm)=\pm2iQ$, with zero cross-pairings. If
$\Phi(0)=0$, $\Phi'(0)=I$, then
$Z=\Phi H_+(0)^{-\dagger}$ obeys $W(H_+,Z)=I$ and $W(Z,Z)=0$.
Invertibility of $H_+(0)$ follows directly from its positive current.
Writing $Z=H_+C_++H_-C_-$ gives
$C_+=(2iQ)^{-1}$ and $C_-^\dagger QC_-=(1/4)Q^{-1}$.
It follows that $\|C_\pm\|_2\leq1/(2q_{\min})$ and therefore
\begin{equation}
 \|Z(r)\|_2\leq\frac{h(r)}{q_{\min}},\qquad
 \|Z'(r)\|_2\leq\frac{q_{\max}}{q_{\min}}+h(r)-1.
\end{equation}
The checked $A(0)<25/4$, $q_{\min}=5/2$, $q_{\max}=23/5$ give
$h(r)<e^{5/2}<13$ and the weights $26/5$, $14$. The rational estimate
$e\leq87/32<11/4$, $\sqrt e<5/3$ yields
$e^{5/2}<605/48<13$; it uses the exponential series remainder after
$n=3$, bounded by $(1/24)/(1-1/5)$.

For the fixed piecewise smooth nominal response $\widehat U$, define
$ I_F=\int_0^R\|F_*-F_n\|_2$, $I_W=\int_0^R\|W_*-W_n\|_2\|\widehat U\|_2$,
$I_R=\int_0^R\|L_n\widehat U-F_n\|_2$,
$J_U=\sum\|[\widehat U]\|_2$, $J_D=\sum\|[\widehat U']\|_2$,
$U_0=\|\widehat U(0)\|_2$ and $U_R=\|\widehat U(R)\|_2$.
Integrating Green's identity on each smooth piece keeps the origin,
jump and exterior boundary terms. Bounding each by the exact adjoint
weights yields
\begin{equation}
 \|C_*-C_{\mathrm{num}}\|_2
 \leq\frac{26}{5}(I_F+I_W+I_R+J_D+B)+14(J_U+U_0)+O,
 \label{supp:complete-budget}
\end{equation}
where
$B=\|b_*\|_2+\epsilon_DU_R+\|b_n\|_2$ and
$O=\epsilon_TU_R+\|c_*\|_2+\|d_{\mathrm{out}}\|_2$.
Here $\epsilon_D,\epsilon_T$ bound exact-to-nominal maps, $b_n$ is
the nominal response's actual boundary mismatch and $d_{\mathrm{out}}$
is its extraction difference from the fixed comparison coefficient.
$W$ denotes the full reduced operator potential, including the
frequency part. Thus frequency errors are retained in $I_W$, and
exterior source terms occur in $b_*,c_*$ exactly once. The archived
component result encloses every term of \eqref{supp:complete-budget};
its positive second-channel margin is the one quoted in the main text.

\subsection{Weighted regularity and the continuation assumption}
The already derived weighted-regularity lemma permits $\gamma=1/10$.
Here is its logical link to the exact data, without a new numerical
Sobolev-norm calculation. From \eqref{supp:exact-tails}, write
$a=0.495$, $b=0.287$. The profile equation gives
$(r^2f')'=r^2N(f)$ with integrable right-hand side; hence
\begin{equation}
 f'(r)=\frac{\ell_\infty}{r^2}
       -r^{-2}\int_r^\infty s^2N(f(s))\,\dd s.
\end{equation}
Integration yields $f(r)=-\ell_\infty/r+O(e^{-ar})$; the known
exponential bound forces $\ell_\infty=0$, so $f'=O(e^{-ar})$.
For the fundamental mode $y=(u,v)$, the bounded exterior matrix in
$y''=M(r)y$ gives $y''=O(e^{-br})$. Its integrability gives a limit
of $y'$, which must be zero because $y\to0$. Thus $y'=O(e^{-br})$.
Twice differentiating these actual ODEs propagates the same rates through
order four. This differentiates equations, not a modulus inequality.
The origin expansions in Section~\ref{supp:matching-positive} give
smooth Cartesian fields; the compact intermediate region is smooth.
Cartesian derivatives through order four are finite sums of radial
derivatives times bounded angular factors and inverse powers of $r$.
After multiplication by $e^{\gamma\sqrt{1+r^2}}$, they remain square
integrable because $a-\gamma>0$, $b-\gamma>0$. Thus the background
and every finite complex combination of its three dipole tangents
belong to $H^4_\gamma(\mathbb R^3)$.

What remains an explicit assumption is that a hypothetical family is
$C^2$ in amplitude in $C^2(\mathbb T;H^4_\gamma)$ with this one fixed
positive weight, smooth real frequencies and the specified first
harmonic tangent (possibly plus a static first tangent).
The nonzero coefficient excludes that class. Pointwise smoothness at
each radius, changing weights along the family, other first harmonics
or circular polarization are not covered by this monopole argument.
